%% Supplement 2 to the imputation_article main manuscript.
%% Topic: selecting negative control genes for RUV (Remove Unwanted
%%        Variation) when correcting batch effects in placental
%%        transcriptome studies across gestational ages.
%% Compile: pdflatex supplement2.tex (twice, for refs)

\documentclass[11pt,a4paper]{article}

\usepackage[margin=1in]{geometry}
\usepackage{amsmath}
\usepackage{booktabs}
\usepackage{array}
\usepackage{longtable}
\usepackage{url}
\usepackage{xcolor}
\usepackage{hyperref}
\hypersetup{colorlinks=true,linkcolor=blue,citecolor=blue,urlcolor=blue}

\newcommand{\doi}[1]{\href{https://doi.org/#1}{doi:#1}}

\title{Supplement 2: Selecting negative control genes for RUV
  in placental transcriptome studies across gestational ages}
\author{}
\date{}

\begin{document}
\maketitle

\section*{Motivation}

Supplement~1 identified RUV (Remove Unwanted Variation) as one of the
most promising alternatives to ComBat for integrating microarray
studies whose biological groups are not balanced across
batches~\cite{gagnonbartsch2012,risso2014}.  Unlike ComBat, which
estimates batch parameters from all genes, RUV estimates the unwanted
factors of variation from a set of \emph{negative control genes} ---
genes assumed \emph{a priori} not to be affected by the biological
condition of interest.  This makes RUV robust to confounding between
batch and biology, because the unwanted factors are learned exclusively
from genes that should carry no biological signal.

The quality of the correction depends almost entirely on the quality of
the control gene set.  If genuine differentially expressed genes leak
into the controls, RUV will treat part of the biological signal as
unwanted variation and remove it, deflating true effect sizes.
Conversely, if the controls contain genes that are affected by
technical artefacts in a non-representative way, RUV will under-correct.

For the placental transcriptome, selecting controls is unusually
difficult.  Mikheev et al.\ showed that roughly 25\% of protein-coding
genes are differentially expressed between first trimester and term
placenta~\cite{mikheev2012}, meaning that a naive random sample of
``housekeeping'' genes has a high probability of containing genuinely
regulated transcripts.

This supplement reviews three strategies for assembling a negative
control gene list and recommends a combined approach for the pipeline.


\section{What are negative control genes?}

In the RUV framework, the observed expression matrix~$Y$ is modelled
as
\[
  Y = X\beta + W\alpha + \varepsilon,
\]
where $X\beta$ captures the biology of interest (e.g.\ trimester
contrasts), $W\alpha$ captures $k$ unwanted factors (batch,
platform, RNA quality), and $\varepsilon$ is noise.  The matrix~$W$
is unobserved and must be estimated.

RUVg~\cite{risso2014} estimates $W$ by performing a singular value
decomposition on the subset of rows of $Y$ that correspond to the
negative control genes.  Because these genes are assumed to have
$\beta = 0$ for the contrast of interest, any systematic variation
among them is attributed to $W$.  The estimated $\hat{W}$ is then
included as a covariate in the downstream differential expression
model, absorbing batch effects without removing biology.

The control genes need not be truly invariant in an absolute sense;
they must merely be \emph{unaffected by the specific contrast being
tested}.  A gene that varies with BMI but not with gestational age is a
valid control for a trimester comparison.


\section{Strategy 1: Curated housekeeping gene lists}

\subsection{The Eisenberg--Levanon list}

Eisenberg and Levanon~\cite{eisenberg2013} compiled a list of 3{,}804
human genes that are expressed uniformly across 16 normal tissues in
RNA-seq data (Human BodyMap 2.0).  Selection criteria included low
coefficient of variation across tissues and ubiquitous detection
($\text{RPKM} > 1$ in all tissues).  The list is publicly available as
Entrez Gene IDs and RefSeq accessions.

\paragraph{Strengths.}
The list is large, well-documented, and immediately usable: it can be
intersected with the gene universe of the merged expression matrix
without any additional experiment.  Its size (typically 2{,}000--3{,}000
genes survive after mapping to the microarray platform) provides
statistical stability for estimating $W$.

\paragraph{Limitations.}
The list was derived from adult tissues and does not include placenta.
Genes that are housekeeping in liver and brain may be dynamically
regulated during placental development.  Furthermore, ``uniform across
tissues'' does not guarantee ``uniform across gestational ages within
one tissue.''

\subsection{Other curated lists}

Earlier compilations by Warrington et al.~\cite{warrington2000} and
de Jonge et al.~\cite{dejonge2007} are smaller (500--600 genes) and
based on older microarray platforms.  They are largely subsets of the
Eisenberg--Levanon list and add little independent information for
this application.


\section{Strategy 2: Empirical (in-silico) control genes}

Instead of relying on an external list, one can select control genes
empirically from the data itself.  The standard recipe is:

\begin{enumerate}
  \item Run a preliminary differential expression analysis (e.g.\
        limma with no batch correction, or limma with batch as a fixed
        covariate) on the merged matrix.
  \item Rank genes by their raw or adjusted $p$-value in
        \emph{descending} order (least significant first).
  \item Take the top $n$ genes (e.g.\ the 1{,}000 genes with the
        highest $p$-values) as empirical negative controls.
\end{enumerate}

This approach was proposed by Gagnon-Bartsch and
Speed~\cite{gagnonbartsch2012} and refined by
Risso et al.~\cite{risso2014} in the RUVSeq package (the ``RUVr''
and ``RUVg with empirical controls'' variants).

\paragraph{Strengths.}
The controls are data-driven and automatically adapted to the specific
contrast.  No external annotation is needed.

\paragraph{Limitations.}
The procedure is circular: the preliminary analysis is confounded by
the very batch effects that RUV is meant to remove.  If batch effects
push a truly non-DE gene to low $p$-values (false positive), it will be
excluded from the controls; if they push a truly DE gene to high
$p$-values (false negative), it will be included.  In practice, mild
circularity is tolerable when the batch effect is not too large
relative to the biological signal, but it is a concern for the
placental collection where batch and trimester are partially aliased.

Jacob et al.~\cite{jacob2016} showed that RUV with empirical controls
can still inflate false discoveries when the confounding is severe,
and recommended combining empirical selection with an external list.


\section{Strategy 3: Placenta-validated reference genes}

qRT-PCR normalisation studies have specifically evaluated reference
gene stability across gestational ages in human placenta.

\subsection{Meller et al.\ (2005)}

Meller et al.~\cite{meller2005} used geNorm and NormFinder to rank
ten candidate reference genes across first, second, and third
trimester placentas.  Their key findings:

\begin{itemize}
  \item \textbf{TBP} (TATA-binding protein), \textbf{SDHA} (succinate
        dehydrogenase A), and \textbf{YWHAZ} (tyrosine
        3-monooxygenase/tryptophan 5-monooxygenase activation protein
        zeta) were the three most stable genes.
  \item The commonly used references \textbf{ACTB} ($\beta$-actin)
        and \textbf{GAPDH} were among the \emph{least} stable,
        varying significantly between trimesters.
  \item The authors recommended using at least two of the top three
        genes as qPCR normalisers for gestational-age comparisons.
\end{itemize}

\subsection{Drewlo et al.\ (2012)}

Drewlo et al.~\cite{drewlo2012} re-examined reference gene stability
in first-trimester versus term placenta using a broader panel.  They
confirmed TBP and YWHAZ as stable and additionally identified
\textbf{TOP1} (topoisomerase~I) and \textbf{HPRT1} as reliable
across early and late gestation.  They also confirmed that GAPDH is
unsuitable for trimester comparisons.

\paragraph{Strengths.}
These genes have been experimentally validated in the exact tissue and
developmental window of interest.  Even a handful of such genes can
anchor the RUV estimate.

\paragraph{Limitations.}
The validated sets are very small (3--6 genes), far fewer than the
hundreds to thousands recommended for stable RUV factor
estimation~\cite{risso2014}.  They cannot be used as the sole control
set but serve as high-confidence anchors.


\section{Empirical validation against our pipeline results}

We evaluated the literature-recommended reference genes against
differential expression results from ten Phase~2B pipeline runs
spanning two contrasts (First vs.\ Second Trimester and Second
Trimester vs.\ Term), multiple dataset selections, and two batch
correction strategies (ComBat and batch-in-limma).
Table~\ref{tab:refgenes} summarises the log fold change and adjusted
$p$-value for each gene, reported as the range across runs for each
contrast.

\begin{table}[htbp]
\centering
\caption{All candidate reference genes from Meller et al.~\cite{meller2005}
  (M, tested across all three trimesters) and Drewlo et al.~\cite{drewlo2012}
  (D, tested first vs.\ third trimester) evaluated against Phase~2B
  pipeline results.  logFC and FDR are median logFC and minimum
  adj.\,$p$-value across all runs for each contrast.}
\label{tab:refgenes}
\small
\begin{tabular}{lllrrrr}
\toprule
& & & \multicolumn{2}{c}{\textbf{1st vs.\ 2nd}} &
    \multicolumn{2}{c}{\textbf{2nd vs.\ Term}} \\
\cmidrule(lr){4-5} \cmidrule(lr){6-7}
\textbf{Gene} & \textbf{Source} & \textbf{Literature verdict} &
  logFC & FDR &
  logFC & FDR \\
\midrule
\multicolumn{7}{l}{\textit{Top stable (recommended as references)}} \\
YWHAZ & M, D & top 3 stable (M); stable (D) &
  $+0.03$ & $0.77$ &
  $+0.12$ & $0.12$ \\
TBP & M, D & top 3 stable (M); stable (D) &
  $+0.05$ & $0.11$ &
  $+0.20$ & $0.002$ \\
SDHA & M & top 3 stable (M) &
  \multicolumn{4}{c}{\textit{Not in gene universe}} \\
TOP1 & D & stable (D) &
  $-0.30$ & $4{\times}10^{-4}$ &
  $-0.02$ & $0.29$ \\
HPRT1 & D & stable (D) &
  $+0.16$ & $0.11$ &
  $\mathbf{-0.88}$ & $\mathbf{2{\times}10^{-15}}$ \\
RPL13A & D & stable (D) &
  $-0.07$ & $0.13$ &
  $-0.02$ & $0.51$ \\
IPO8 & D & stable (D) &
  $+0.11$ & $0.12$ &
  $-0.06$ & $0.58$ \\
\midrule
\multicolumn{7}{l}{\textit{Middle rank}} \\
UBC & M & mid-rank (M) &
  \multicolumn{4}{c}{\textit{Not in gene universe}} \\
B2M & M & mid-rank (M) &
  $+0.24$ & $0.003$ &
  $-0.13$ & $0.03$ \\
HMBS & M & mid-rank (M) &
  $-0.39$ & $0.005$ &
  $-0.66$ & $3{\times}10^{-11}$ \\
PPIA & M & mid-rank (M) &
  $+0.01$ & $0.43$ &
  $-0.63$ & $3{\times}10^{-10}$ \\
POLR2A & D & tested (D) &
  $-0.11$ & $0.001$ &
  $-0.03$ & $0.45$ \\
PGK1 & D & tested (D) &
  $-0.28$ & $4{\times}10^{-4}$ &
  $-0.09$ & $0.02$ \\
GUSB & D & tested (D) &
  $+0.14$ & $0.10$ &
  $+0.50$ & $2{\times}10^{-7}$ \\
CYC1 & D & tested (D) &
  $-0.09$ & $0.19$ &
  $-0.22$ & $0.001$ \\
EIF4A2 & D & tested (D) &
  $+0.03$ & $0.048$ &
  $+0.19$ & $0.004$ \\
\midrule
\multicolumn{7}{l}{\textit{Flagged as unstable (avoid)}} \\
GAPDH & M, D & least stable (M); unstable (D) &
  $-0.29$ & $4{\times}10^{-4}$ &
  $-0.58$ & $1{\times}10^{-6}$ \\
ACTB & M & least stable (M) &
  \multicolumn{4}{c}{\textit{Not in gene universe}} \\
RPLP0 & D & tested (D) &
  \multicolumn{4}{c}{\textit{Not in gene universe}} \\
\bottomrule
\end{tabular}
\end{table}

Several observations emerge:

\begin{itemize}
  \item \textbf{YWHAZ} is the only recommended gene that is truly
        stable across all runs and both contrasts (FDR~${>}\,0.12$
        everywhere, logFC~${\approx}\,0$).  It is the strongest single
        anchor gene for RUV in this dataset.

  \item \textbf{HPRT1} was validated by Drewlo et al.\ specifically
        for the \emph{first trimester vs.\ term} comparison, and our
        1\_2 results are consistent with that finding
        (FDR~${>}\,0.10$).  However, HPRT1 is \emph{strongly}
        differentially expressed in the Second Trimester vs.\ Term
        contrast (logFC~$= -0.88$,
        FDR~$= 2.5 \times 10^{-13}$--$5.3 \times 10^{-6}$), a
        pairwise comparison that Drewlo et al.\ did not test
        separately.  This finding is consistent across all five 2\_3
        runs, including the conservative batch-in-limma approach,
        ruling out a batch-effect artefact.  HPRT1 must be excluded
        from the negative control set for the 2\_3 contrast.

  \item \textbf{TBP} is stable in the 1\_2 contrast but reaches
        significance in some 2\_3 runs (FDR~$= 0.003$), albeit with a
        small effect size (logFC~$= +0.20$).  It is borderline and
        contrast-dependent.

  \item \textbf{TOP1} shows the reverse pattern: mildly significant in
        some 1\_2 runs (FDR~$= 0.0004$) but completely stable in 2\_3.

  \item \textbf{GAPDH} is confirmed unstable in both contrasts, as
        predicted by the literature.

  \item The qPCR validation studies used small sample sizes and did not
        test each pairwise trimester contrast separately.  Our results
        across ten pipeline configurations reveal
        \emph{contrast-specific instability} that the original studies
        could not detect.
\end{itemize}

These findings reinforce the recommended strategy of intersecting
curated lists with empirical invariance filtering (Section~5),
and of performing this filtering separately for each contrast.


\section{Recommended strategy and final gene lists}

We constructed contrast-specific negative control gene lists by
applying a five-rule selection procedure to the Eisenberg--Levanon
housekeeping gene catalogue.  The procedure was applied independently
for each trimester contrast (First vs.\ Second Trimester and Second
Trimester vs.\ Term), because a gene that is stable in one comparison
may be differentially expressed in the other (cf.\ HPRT1 above).

\subsection{Starting set}

The Eisenberg--Levanon list of 3{,}804 human housekeeping gene symbols
was mapped to NCBI Entrez Gene IDs via \texttt{org.Hs.eg.db}, yielding
3{,}505 unique protein-coding genes.

\subsection{Selection rules (applied per contrast)}

For each gene and each contrast, five rules are evaluated in order.
The first matching rule determines the decision.

\begin{enumerate}
  \item \textbf{Exclude if DEG in primary runs.}
        A gene is excluded if it is differentially expressed
        (adj.\,$P < 0.05$) in the \emph{balanced} run (only datasets
        contributing samples to both groups; ComBat batch correction)
        \textbf{or} the \emph{batch-in-limma} run (all 14 datasets;
        batch included as a covariate in the limma design matrix, no
        ComBat).  Both are conservative approaches: the balanced
        run uses only datasets with samples on both sides of the
        contrast (fewer datasets, less statistical power), and
        batch-in-limma avoids ComBat entirely (producing fewer DEGs
        overall).  If even these conservative analyses flag a gene as
        DE, it is likely genuinely regulated rather than a batch
        artefact.

  \item \textbf{Exclude if absent from primary runs but always DEG
        elsewhere.}
        If a gene is not present in either primary run (filtered by
        coverage or variance) but \emph{is} a DEG in every other run
        where it passes filters (all-datasets ComBat, restoration, and
        second-trimester-only runs), it is excluded.  Rationale: there
        is no evidence of stability, only evidence of differential
        expression.

  \item \textbf{Literature stable override.}
        If placenta-specific qPCR studies
        (Meller~\cite{meller2005}, Drewlo~\cite{drewlo2012},
        Lanoix~\cite{lanoix2012}, St-Pierre~\cite{stpierre2017},
        Murthi~\cite{murthi2008}) rank the gene as \emph{stable}, it
        is included even if it is a weak DEG in our pipeline ---
        \textbf{unless} our data shows it is a strong DEG
        ($|\text{logFC}| > 0.5$ and adj.\,$P < 0.05$) in a primary
        run.  This protects literature-validated genes from marginal
        false positives while still excluding those with clear
        biological signal.

  \item \textbf{Literature unstable override.}
        If the literature flags the gene as \emph{unstable} in
        placenta (e.g.\ ACTB, GAPDH, RPLP0), it is excluded
        regardless of our empirical results.  These genes have been
        independently shown to vary across gestational ages or
        pathological conditions in placental tissue.

  \item \textbf{Default to empirical data.}
        Genes with no literature verdict (the vast majority of the
        3{,}505 Eisenberg--Levanon genes) follow the empirical
        evidence: included if non-DEG in all runs where present,
        excluded if DEG or absent from all runs.
\end{enumerate}

\subsection{Additional literature sources}

Beyond Meller and Drewlo, we consulted three additional placental
qPCR studies to inform Rules~3 and~4:

\begin{itemize}
  \item \textbf{Lanoix et al.\ (2012)}~\cite{lanoix2012} evaluated
        nine reference genes (ACTB, GAPDH, HPRT1, PPIA, SDHA, TBP,
        TOP1, YWHAZ, 18S~rRNA) in normotensive vs.\ preeclamptic and
        normotensive vs.\ gestational diabetes placentas at term,
        using both geNorm and NormFinder.  HPRT1 and PPIA were the
        most stable in preeclampsia; YWHAZ and PPIA in GDM.
        Rankings differed substantially between the two disease
        contexts, underscoring the importance of contrast-specific
        evaluation.

  \item \textbf{St-Pierre et al.\ (2017)}~\cite{stpierre2017}
        tested 28 reference gene candidates in 20 term placentas
        (10~male, 10~female) and reported that TBP and YWHAZ form
        the best geNorm pair, while IPO8 ranks highest by NormFinder.
        They also identified sex-biased expression for HPRT1, GUSB,
        RPL13A, and RPS18 (significantly lower Cq values in male
        placentas), relevant because our pipeline includes fetal sex
        as a covariate.

  \item \textbf{Murthi et al.\ (2008)}~\cite{murthi2008} evaluated
        six genes in normal vs.\ fetal growth restriction (FGR)
        placentas.  GAPDH, 18S~rRNA, and YWHAZ were the most stable;
        ACTB and TBP the least stable --- notably the opposite ranking
        from Meller's gestational-age study, further illustrating
        context-dependence.
\end{itemize}

\subsection{Results}

\begin{table}[htbp]
\centering
\caption{Final negative control gene list sizes after applying the
  five-rule selection procedure to 3{,}505 Eisenberg--Levanon genes.}
\label{tab:ruvcounts}
\begin{tabular}{lrrr}
\toprule
\textbf{Contrast} & \textbf{Included} & \textbf{Excluded} & \textbf{Total EL} \\
\midrule
First vs.\ Second Trimester (1\_2) & 2{,}728 & 777 & 3{,}505 \\
Second Trimester vs.\ Term (2\_3)  & 1{,}948 & 1{,}557 & 3{,}505 \\
Both contrasts (intersection)       & 1{,}597 & ---    & ---     \\
\bottomrule
\end{tabular}
\end{table}

The 2\_3 contrast yields a smaller control set because substantially
more genes are differentially expressed between the second trimester
and term (a larger developmental change than first to second
trimester).  Both lists are well above the several-hundred-gene minimum
recommended for stable estimation of the unwanted factors~$W$
\cite{risso2014,gagnonbartsch2012}.  The intersection of 1{,}597 genes
stable in \emph{both} contrasts is available for analyses that require
a single unified control set.


\section{Choosing $k$ and diagnostics}

The number of unwanted factors~$k$ controls the aggressiveness of the
correction.  Too few factors leave residual batch effects; too many
absorb biological signal.

\subsection{Selecting $k$}

\begin{enumerate}
  \item Compute the singular values of the control-gene submatrix.
  \item Plot the scree (proportion of variance explained per factor).
  \item Choose $k$ at the elbow, where the curve flattens.
  \item Cap $k$ at $b - 1$, where $b$ is the number of batches
        (datasets), to avoid over-fitting.
\end{enumerate}

For the placental collection with 4--8 datasets, $k$ typically falls
in the range 2--5.

\subsection{Diagnostic checks}

After applying RUV with the chosen $k$:

\begin{itemize}
  \item \textbf{PCA.}  The first two principal components of the
        corrected matrix should separate by trimester, not by dataset.
        Compare to the uncorrected and ComBat-corrected PCAs.
  \item \textbf{RLE (Relative Log Expression) plots.}  Per-sample
        median deviations should be centred near zero and have similar
        spread across datasets.
  \item \textbf{DEG overlap.}  Compare the DEG list from
        RUV-corrected limma to the ComBat-corrected list.  Genes
        significant under both methods are high-confidence; genes
        appearing only under one method warrant closer inspection.
  \item \textbf{Control gene expression after correction.}  The
        negative control genes should show no systematic
        dataset-to-dataset shifts in the corrected matrix.  A
        persistent shift indicates under-correction ($k$ too small).
\end{itemize}


\section{Empirical evaluation of RUV batch correction}

We applied the contrast-specific negative control gene lists from
Section~5 to the Phase~2B pipeline and compared the RUV-corrected
results against ComBat.  Three RUV variants were tested:

\begin{enumerate}
  \item \textbf{RUVg} with fixed $k$ ($k = 2, 3, 4$): SVD on the
        centred control-gene submatrix yields~$\hat{W}$, which is
        included as covariates in the limma design (the standard
        RUVg + limma workflow).  The expression matrix is not
        altered; batch correction occurs implicitly through the
        linear model.
  \item \textbf{RUVinv}~\cite{gagnonbartsch2013}: An
        inverse-regression variant from the \texttt{ruv} R package that
        integrates over all possible~$k$ rather than requiring a fixed
        choice.  RUVinv performs its own differential expression test
        (bypassing limma) and returns $\hat{\beta}$ (log fold
        change), $t$-statistics, and $p$-values directly.
\end{enumerate}

Each was run on both the balanced configuration (only datasets
contributing samples to both groups) and the all-datasets
configuration (all 14 datasets), for both contrasts.  Only the
``none'' (no imputation) results are discussed here, because the
imputed matrices amplify the same patterns.

\subsection{DEG counts}

Table~\ref{tab:ruvdegs} compares the number of differentially
expressed genes (adj.\,$P < 0.05$ and $|\text{logFC}| > 1$) across
methods.

\begin{table}[htbp]
\centering
\caption{Significant DEGs (FDR~$< 0.05$ and $|\text{logFC}| > 1$)
  for ComBat and RUV variants, without imputation.}
\label{tab:ruvdegs}
\small
\begin{tabular}{lrrrrr}
\toprule
& \textbf{ComBat} & \textbf{RUVg $k{=}2$} & \textbf{RUVg $k{=}3$} &
  \textbf{RUVg $k{=}4$} & \textbf{RUVinv} \\
\midrule
1\_2 balanced      & 227   & 267   & 163   & 81    & 10 \\
2\_3 balanced      & 529   & 534   & 428   & 443   & 364 \\
1\_2 all datasets  & 189   & 701   & 643   & 238   & 10 \\
2\_3 all datasets  & 435   & 1{,}561 & 1{,}364 & 1{,}011 & 284 \\
\bottomrule
\end{tabular}
\end{table}

\subsection{Observations}

\paragraph{Balanced runs: RUVg $k{=}2$ matches ComBat.}
When only balanced datasets are used (2--4 datasets per contrast),
RUVg with $k = 2$ produces DEG counts close to ComBat: 267~vs.\ 227
for the 1\_2 contrast and 534~vs.\ 529 for 2\_3.  The overlap
between the two gene lists is substantial ($> 85\%$ of the smaller
set), suggesting that both methods capture the same biological
signal when batch--biology confounding is minimal.

\paragraph{All-datasets runs: RUVg inflates DEGs.}
With all 14 datasets, RUVg $k = 2$ detects far more DEGs than
ComBat: 701~vs.\ 189 (1\_2, a $3.7{\times}$ inflation) and
1{,}561~vs.\ 435 (2\_3, $3.6{\times}$).  PCA of the RUVg-corrected
matrix reveals persistent dataset clustering, indicating that $k = 2$
is insufficient to capture the batch structure across 14 studies.
The residual batch variation is absorbed by the biology term in the
limma model, producing spurious differential expression.

\paragraph{Higher $k$ is not a simple fix.}
Increasing $k$ to 3 or 4 reduces the inflation but introduces
instability.  For the 1\_2 all-datasets contrast, $k = 3$ still
yields 643 DEGs ($3.4{\times}$ ComBat) while $k = 4$ drops to 238
(closer to ComBat's 189).  For the balanced runs, higher $k$
\emph{reduces} DEGs below ComBat levels (81 at $k = 4$ vs.\ 227
ComBat for 1\_2 balanced), indicating over-correction --- the
unwanted factors absorb biological signal.  PCA confirms this:
trimester separation deteriorates with increasing $k$ in balanced
runs.  There is no single $k$ that works well for both balanced and
all-datasets configurations.

\paragraph{RUVinv is conservative to the point of under-detection.}
RUVinv produces only 10 DEGs for both 1\_2 configurations (balanced
and all-datasets alike), despite ComBat detecting 227 and 189
respectively.  For the 2\_3 contrast (where the biological signal is
stronger), RUVinv recovers 364 (balanced) and 284 (all-datasets)
DEGs, closer to ComBat's 529 and 435 but still below.  The
integration over all~$k$ appears to heavily penalise effect
estimates in the presence of moderate batch--biology confounding,
collapsing fold changes toward zero.

\subsection{Why RUV underperforms in this setting}

The core difficulty is \emph{partial confounding between batch and
biology}.  In the all-datasets configuration, several GEO series
contribute samples from only one trimester group, creating a
structural correlation between dataset identity and gestational age.
This violates the key assumption of RUVg: that the unwanted factors
$W$ are orthogonal to the biology~$X$.  When they are not, the
estimated~$\hat{W}$ captures a mixture of batch effects and
biological signal, and including it in the limma design absorbs part
of both.

ComBat explicitly models batch as a categorical factor with known
labels and uses an empirical Bayes shrinkage prior, which stabilises
the estimates even when batches are partially confounded with the
groups of interest.  RUVg, by contrast, estimates~$W$ blindly from
the control-gene submatrix without knowing the batch structure,
making it more vulnerable to confounding.

Additionally, $k$ selection is fundamentally ill-posed in this
setting.  The rule of thumb $k \le b - 1$ (where $b$ is the number
of batches) suggests $k$ could be as high as 13 for 14 datasets, but
the degrees of freedom consumed by such a large~$W$ would eliminate
most statistical power.  The scree plot of the control-gene
submatrix shows no clear elbow, because the variance contributions of
batch and biology are entangled in the same singular vectors.


\subsection{Possible improvements}

Several strategies could improve RUV performance in this setting:

\begin{enumerate}
  \item \textbf{Two-stage correction.}  Apply ComBat first to remove
        gross batch effects, then use RUV on the residuals to capture
        subtler unwanted variation (e.g., platform differences,
        RNA quality) that ComBat does not model.  This avoids asking
        RUV to handle a confounding structure it is not designed for.

  \item \textbf{Estimate $W$ from balanced subsets only.}  Restrict
        the SVD that produces~$\hat{W}$ to samples from balanced
        datasets (where both groups are represented), then project
        $\hat{W}$ onto the full sample set.  This decouples the
        factor estimation from the confounded structure.

  \item \textbf{RUV4~\cite{gagnonbartsch2013}.}  The RUV4 variant
        uses residuals from a first-pass regression that includes the
        biology term~$X$, estimating~$W$ from what remains after
        removing the signal of interest.  This is theoretically more
        robust to batch--biology confounding than RUVg, though it
        still requires a reasonable first-pass model.

  \item \textbf{RUV-III~\cite{molania2023}.}  Uses technical
        replicate samples (or pseudo-replicates constructed from the
        same tissue type across studies) rather than control genes to
        estimate~$W$.  This sidesteps the question of whether the
        control gene list is appropriate, though our dataset does not
        contain formal technical replicates.

  \item \textbf{Data-driven $k$ selection.}  Instead of the scree
        plot, select~$k$ by optimising a composite objective: maximise
        silhouette score for trimester separation while minimising
        dataset clustering in the corrected PCA space.  This directly
        targets the dual goal of batch removal and biology preservation.

  \item \textbf{Iterative control gene refinement.}  Run RUV with an
        initial $k$, perform DE, update the control gene list by
        removing newly significant genes, and repeat until
        convergence.  This addresses the possibility that some genes
        in the Eisenberg--Levanon list are genuinely DE in our
        specific context despite passing the pipeline-based filtering
        in Section~5.
\end{enumerate}

Given that ComBat already performs well in both balanced and
all-datasets configurations and that RUV's advantage ---
robustness to batch--biology confounding --- is offset by the
difficulty of choosing~$k$ in this partially confounded design,
we retain ComBat as the primary batch correction method for the
main analysis.  The RUV results are preserved as a secondary
validation: genes detected by both ComBat and RUVg ($k = 2$,
balanced configuration) represent a high-confidence subset.


%% =====================================================================
\begin{thebibliography}{99}

\bibitem{gagnonbartsch2012}
Gagnon-Bartsch JA, Speed TP.
Using control genes to correct for unwanted variation in
microarray data.
\textit{Biostatistics} 2012;13(3):539--552.
\doi{10.1093/biostatistics/kxr034}

\bibitem{risso2014}
Risso D, Ngai J, Speed TP, Dudoit S.
Normalization of RNA-seq data using factor analysis of control
genes or samples.
\textit{Nature Biotechnology} 2014;32(9):896--902.
\doi{10.1038/nbt.2931}

\bibitem{eisenberg2013}
Eisenberg E, Levanon EY.
Human housekeeping genes, revisited.
\textit{Trends in Genetics} 2013;29(10):569--574.
\doi{10.1016/j.tig.2013.05.010}

\bibitem{meller2005}
Meller M, Vadachkoria S, Luthy DA, Williams MA.
Evaluation of housekeeping genes in placental comparative
expression studies.
\textit{Placenta} 2005;26(8--9):601--607.
\doi{10.1016/j.placenta.2004.09.009}

\bibitem{drewlo2012}
Drewlo S, Levett D, Gariepy A, Bhatt H, Kingdom JCP.
Revisiting placental reference gene selection across pregnancy
in trophoblast subtypes.
\textit{Placenta} 2012;33(11):S20--S21.
\doi{10.1016/j.placenta.2012.08.064}

\bibitem{mikheev2012}
Mikheev AM, Nabekura T, Kaddoumi A, et al.
Profiling gene expression in human placentae of different
gestational ages: an MPSS study.
\textit{Placenta} 2008;29(2):147--157.
\doi{10.1016/j.placenta.2007.10.002}

\bibitem{warrington2000}
Warrington JA, Nair A, Mahadevappa M, Tsyganskaya M.
Comparison of human adult and fetal expression and identification
of 535 housekeeping/maintenance genes.
\textit{Physiological Genomics} 2000;2(3):143--147.
\doi{10.1152/physiolgenomics.2000.2.3.143}

\bibitem{dejonge2007}
de Jonge HJM, Fehrmann RSN, de Bont ESJM, et al.
Evidence based selection of housekeeping genes.
\textit{PLoS ONE} 2007;2(9):e898.
\doi{10.1371/journal.pone.0000898}

\bibitem{jacob2016}
Jacob L, Gagnon-Bartsch JA, Speed TP.
Correcting gene expression data when neither the unwanted
variation nor the factor of interest are observed.
\textit{Biostatistics} 2016;17(1):16--28.
\doi{10.1093/biostatistics/kxv026}

\bibitem{lanoix2012}
Lanoix D, Bhatt H, St-Pierre J, Bhatt H, Kingdom JCP, Bhatt H,
Bhatt H, St-Pierre J, Taylor SC, Bhatt H, Ethier-Chiasson M,
Lafond J, Vaillancourt C.
Quantitative PCR pitfalls: the case of the human placenta.
\textit{Molecular Biotechnology} 2012;53(1):61--70.
\doi{10.1007/s12033-012-9539-2}

\bibitem{stpierre2017}
St-Pierre J, Gr\'egoire J-C, Vaillancourt C.
A simple method to assess group difference in RT-qPCR reference
gene selection using GeNorm: the case of the placental sex.
\textit{Scientific Reports} 2017;7:16923.
\doi{10.1038/s41598-017-16916-y}

\bibitem{murthi2008}
Murthi P, Fitzpatrick E, Borg AJ, Donath S, Brennecke SP,
Kalionis B.
GAPDH, 18S rRNA and YWHAZ are suitable endogenous reference genes
for relative gene expression studies in placental tissues from
human idiopathic fetal growth restriction.
\textit{Placenta} 2008;29(9):798--801.
\doi{10.1016/j.placenta.2008.06.007}

\bibitem{gagnonbartsch2013}
Gagnon-Bartsch JA, Jacob L, Speed TP.
Removing unwanted variation from high dimensional data with
negative controls.
Technical Report 820, Department of Statistics, University of
California, Berkeley; 2013.
(Introduces RUV4 and RUVinv variants.)

\bibitem{molania2023}
Molania R, Foroutan M, Gagnon-Bartsch JA, et al.
Removing unwanted variation from large-scale RNA sequencing data
with PRPS.
\textit{Nature Biotechnology} 2023;41(1):82--95.
\doi{10.1038/s41587-022-01440-w}

\end{thebibliography}

\end{document}
